\needspace{10\baselineskip}
\section{进程线程调度}

调度要解决三个问题 \bk：\textbf{What}——按什么原则挑选进程/线程（调度算法）；\textbf{When}——什么时候分配处理器（调度时机）；\textbf{How}——如何分配 CPU，即进程的上下文切换（调度过程）。其中上下文切换是\textbf{机制}，调度算法是\textbf{策略}。老师的答题要求：调度算法\textbf{不要背优缺点，要像提方案一样多维度分析}——是否抢占？会不会饥饿？对 I/O 密集型/CPU 密集型、长/短进程是否友好？开销多大？响应时间、周转时间、吞吐量如何？是否需要预知运行时间？适用于什么场景？

\subsection{重点问题}

\Q{Q4-1.1 为什么操作系统需要“调度”？CPU 只有一个时，怎么同时跑多个程序？}
\textbf{需求多于资源}：系统中有 N 个进程处于就绪状态，等待在 M 个 CPU 上运行（M $\ge$ 1，通常 M 远小于 N），必须决定“把哪个 CPU 分给哪个进程”。调度就是\textbf{控制、协调进程对 CPU 的竞争}：按一定的调度算法从就绪队列中选择一个进程，把 CPU 的使用权交给它。完成这一选择的内核函数称为\textbf{调度程序}。另一个动机是\textbf{提高资源利用率}：一个进程等 I/O 时，把 CPU 交给别的进程，这正是多道程序设计的出发点。

只有一个 CPU 时，靠\textbf{分时复用}“同时”运行多个程序：CPU 时间被切成毫秒级的时间片，轮流分配给各个进程，切换得足够快，人就感觉不到停顿。要注意这是\textbf{并发}（宏观上同时、微观上交替），而不是并行。

\Q{Q4-1.2 调度这件事什么时候做？做的理由有哪些？}
调度时机详见 4.2 节，归纳为下表：\par
\begin{btab}{P{0.300}P{0.700}}
\thead{时机 & 为什么要调度}
\trow{进程退出、出错终止 & CPU 空出来了，必须选一个新进程}
\trow{进程阻塞（等 I/O、等资源、\cd{wait}） & 同上，CPU 空出来了}
\trow{进程主动让出（\cd{yield}） & 进程自己要求把 CPU 让给别人}
\trow{时间片用完（时钟中断） & 保证公平和响应时间，不让一个进程长期霸占 CPU}
\trow{新进程创建、阻塞进程被唤醒 & 就绪队列中多了一个进程，它可能比当前进程优先级更高，需要判断是否抢占}
\end{btab}
\textbf{实现上的统一时刻}：\imp{内核处理完中断/异常/系统调用、返回用户态之前的最后时刻}检查是否需要调度。事件处理过程中只是做一个标记（Linux 设置当前进程的 \cd{TIF_NEED_RESCHED}），到统一的出口再真正调用调度程序；xv6 中对应 \cd{usertrap} 末尾的 \cd{if(which_dev == 2) yield();}。

\Q{Q4-1.3 如果没有可被调度的进程，系统做什么？（什么是 idle 进程？）}
运行 \textbf{idle 进程}（系统空闲进程）。它的优先级最低，并且永远处于就绪状态，这样调度程序总能选出一个进程，不必为“就绪队列为空”写特殊的处理逻辑。idle 进程\textbf{能被任何事件打断}：一旦有中断使某个进程就绪，就立即切换过去。它通常执行省电的等待指令（x86 的 \cd{hlt}，RISC-V 的 \cd{wfi}），让 CPU 进入低功耗状态直到下一个中断。

各系统的实现：Linux 每个 CPU 有一个 idle 任务（PID 为 0，名为 \cd{swapper/N}），配合 NO\_HZ 在空闲时停掉周期时钟中断以省电；Windows 每个 CPU 有一个空闲线程，优先级为 0，它还会处理 DPC（延迟过程调用），零页线程则利用空闲时间把空闲页清零；xv6 没有专门的 idle 进程，\cd{scheduler()} 扫描一遍进程表没找到 RUNNABLE 的进程，就执行 \cd{wfi} 等待中断。\cd{top} 命令中的 \texttt{\%id} 就是 CPU 处于空闲状态的时间比例。

\Q{Q4-1.4 上下文切换的过程？有哪些开销？}
\textbf{过程}（场景：进程 A 下 CPU，进程 B 上 CPU）：
\begin{enumerate}
    \item 保存进程 A 的上下文环境（程序计数器、程序状态字、栈指针、其他寄存器）；
    \item 用新的状态和其他相关信息更新 A 的 PCB（如从运行改为就绪或阻塞）；
    \item 把 A 移到合适的队列（就绪队列或某个等待队列）；
    \item 把 B 的状态设为运行；
    \item \textbf{切换地址空间}：把 B 的页表基址装入页表基址寄存器（x86 CR3，RISC-V \cd{satp}）；
    \item 从 B 的 PCB 中恢复其上下文环境，B 从上次被打断的地方继续执行。
\end{enumerate}
实际系统中，在保存上下文之前先检查 \cd{need_resched}：如果不需要切换，就直接从内核栈上恢复原进程的寄存器返回，不做后面的步骤。

\textbf{开销}分两类：
\begin{itemize}
    \item \textbf{直接开销}：内核执行切换代码本身花费的 CPU 时间——保存和恢复寄存器、执行调度算法、切换地址空间（写 \cd{satp} 并执行 \cd{sfence.vma} 刷新 TLB）。
    \item \textbf{间接开销}：切换之后，新进程的数据不在 cache 中、地址映射不在 TLB 中、文件块可能不在 buffer cache 中，开始运行后会产生大量 miss，降低速度。这部分往往比直接开销更大。
\end{itemize}
\textbf{缓解方法}：TLB 表项带上进程标识（x86 的 PCID，RISC-V 的 \textbf{ASID}），切换时不必刷新整个 TLB；同一进程的线程之间切换不需要换页表；协程切换不进内核；亲和性调度让进程尽量回到原来的 CPU，利用残留的 cache。\textbf{观察方法}：\cd{vmstat 1} 的 cs 列是每秒上下文切换次数；\cd{/proc/PID/status} 中的 \cd{voluntary_ctxt_switches}（自愿切换，I/O 密集型进程多）和 \cd{nonvoluntary_ctxt_switches}（被抢占，CPU 密集型进程多）。

\Q{Q4-1.5 关于调度算法，我们都关心什么？}
关心一组性能指标（4.3 节）：从用户角度看是周转时间、响应时间、截止时间、可预测性；从系统角度看是吞吐量、CPU 利用率、公平性、均衡性；特定场景还要考虑实时性和能耗。

\Q{Q4-1.6 不同类型的操作系统都适用同一种调度算法吗？}
不适用，因为不同系统的追求目标不同：
\begin{itemize}
    \item \textbf{批处理系统}：追求吞吐量大、周转时间短、CPU 利用率高。适用 FCFS、SJF、SRTN、HRRN。
    \item \textbf{交互式系统}：追求响应时间短（一般要求在 50–150ms 以内）、公平、均衡。适用 RR、优先级调度、多级反馈队列、CFS 等。
    \item \textbf{实时系统}：追求满足截止时间、行为可预测。适用 RM、EDF、固定优先级抢占调度。
\end{itemize}
现代通用 OS 是\textbf{混合}的：把进程分成不同类别，多种策略并存。例如 Linux 中实时进程用 \cd{SCHED_FIFO}/\cd{SCHED_RR}，普通进程用 CFS/EEVDF。

\Q{Q4-1.7 调度算法的目标到底是什么？是越快越好吗？}
\textbf{不是越快越好}，而是在\textbf{相互冲突的目标之间折中}：
\begin{itemize}
    \item 响应时间短 vs 吞吐量大：时间片越短，响应越快，但切换越频繁，用于有效工作的时间越少；
    \item 平均周转时间短 vs 公平：SJF 平均周转时间最短，但长作业可能被饿死；
    \item 优先级 vs 公平：照顾重要的进程，就要牺牲其他进程；
    \item 平均性能 vs 可预测性：实时系统宁可平均慢一点，也要保证最坏情况；
    \item 性能 vs 能耗：移动设备上要权衡。
\end{itemize}
类比金融中投资组合的“收益—风险”有效前沿：没有一个点在所有维度上都最好，只能根据系统的目标选择合适的折中点。

\Q{Q4-1.8 选进程时都考虑了哪些因素？单一因素还是多因素？}
早期算法只考虑\textbf{单一因素}：FCFS 只看到达顺序，SJF 只看运行时间，最高优先级只看优先级。好的算法往往\textbf{综合多个因素}：HRRN 同时考虑等待时间和运行时间；MLFQ 考虑优先级、时间片和进程的历史行为；Windows 考虑基本优先级、动态提升和时间配额；Linux O(1) 考虑 nice 值和平均睡眠时间；CFS 考虑权重和实际运行时间（另外还有唤醒补偿、亲和性、负载均衡）。还有一些隐含的因素：是 I/O 密集型还是 CPU 密集型、亲和性（上次在哪个 CPU 上运行）、属于哪个用户或 cgroup、是否持有锁。

\Q{Q4-2.1 适用于批处理系统的调度算法有哪些？}
\textbf{先来先服务 FCFS}、\textbf{短作业优先 SJF}、\textbf{最短剩余时间优先 SRTN}、\textbf{最高响应比优先 HRRN}。它们关注吞吐量和周转时间，详见 4.4 节。

\Q{Q4-2.2 适用于交互式系统的调度算法有哪些？}
\textbf{时间片轮转 RR}（及其改进虚拟轮转 VRR）、\textbf{最高优先级 HPF}、\textbf{多级队列}、\textbf{多级反馈队列 MLFQ}、\textbf{最短进程优先 SPN}，以及按比例分配的\textbf{公平共享、保证调度、彩票调度}。它们关注响应时间和公平性，详见 4.5 节。

\Q{Q4-2.3 适用于实时系统的调度算法有哪些？}
静态优先级的\textbf{速率单调调度 RM} 和动态优先级的\textbf{最早截止时间优先 EDF}，详见 4.9 节。

\Q{Q4-2.4 什么是抢占式调度？什么是非抢占式调度？区别在哪里？}
\begin{itemize}
    \item \textbf{非抢占式}（不可剥夺）：调度程序一旦把 CPU 分配给某个进程，就让它一直运行下去，直到它\textbf{自己}放弃 CPU（结束、因等待事件而阻塞、主动 \cd{yield}）为止。
    \item \textbf{抢占式}（可剥夺）：进程正在运行时，OS 可以\textbf{强行}剥夺它的 CPU，分配给其他进程。抢占的原因通常是有更高优先级的进程就绪，或时间片用完。
\end{itemize}
“抢占”是可以与其他规则组合的“\textbf{前缀属性}”：SJF + 抢占 = SRTN，优先级调度可以是抢占式或非抢占式。\textbf{区别}：
\begin{itemize}
    \item \textbf{调度时机不同}：非抢占式只在进程退出、阻塞、主动让出时调度；抢占式还要在时钟中断（时间片到）、高优先级进程就绪时考虑调度。
    \item \textbf{硬件要求不同}：\imp{抢占式调度必须有时钟中断}，否则 OS 一旦把 CPU 交出去就没有机会再拿回来。
    \item \textbf{健壮性不同}：非抢占式下，一个陷入死循环的进程会卡死整个系统；抢占式可以防止一个进程霸占 CPU，保证响应时间。
    \item \textbf{代价不同}：抢占式切换更频繁、开销更大；进程可能在任意时刻被打断，共享数据必须加锁保护，并发问题更多。
\end{itemize}
例：FCFS、SJF、Windows 3.x 的协作式多任务、协程是非抢占式的；RR、SRTN、Linux、Windows NT 以后是抢占式的。删掉 xv6 \cd{usertrap} 中的 \cd{if(which_dev == 2) yield();}，它就变成非抢占式调度了。另外还要区分用户态抢占和\textbf{内核抢占}：Linux 2.4 的内核不可抢占（进程在内核态运行时不会被切换），2.6 起可以配置 \cd{CONFIG_PREEMPT} 支持内核抢占。

\Q{Q4-2.5 从哪几方面对调度算法进行比较？}
从以下维度比较（见 4.6 节对比表）：\textbf{选择函数}（依据什么选进程）、\textbf{决策模式}（抢占或非抢占、何时抢占）、\textbf{吞吐量}、\textbf{响应时间}、\textbf{开销}、\textbf{对进程的影响}（对短进程、长进程、I/O 密集型进程是否有利）、\textbf{是否会饥饿}。再加上公平性和是否需要预知运行时间。

\Q{Q4-2.6 机制与策略分离的原则在调度中的应用？}
\textbf{为什么要分离}：调度策略需要随着应用场景变化而调整，而切换进程这套底层代码复杂、与硬件相关，应当保持稳定。分离后，换一种调度策略不必改动底层机制。\textbf{机制}：上下文切换（\cd{swtch}）、就绪队列等数据结构、时钟中断、\cd{schedule()} 的框架；\textbf{策略}：选谁、时间片多长、何时抢占。\textbf{怎么做}：
\begin{enumerate}
    \item \textbf{内核模块化}：Linux 用 \cd{struct sched_class} 把调度算法做成可插拔的“调度类”，\cd{schedule()} 依次询问 stop $\to$ deadline（EDF）$\to$ rt（FIFO/RR）$\to$ fair（CFS/EEVDF）$\to$ idle 各调度类，找到第一个有就绪任务的类；新的 sched\_ext 甚至允许用 eBPF 在运行时加载调度策略。
    \item \textbf{调度机制在内核中，参数由用户进程设置}：父进程最了解它的哪个子进程或线程更重要，所以内核提供 \cd{nice}、\cd{setpriority}、\cd{sched_setscheduler}、\cd{sched_setaffinity} 等系统调用，让用户进程设置调度参数。
\end{enumerate}
“按优先级选进程”是机制；优先级继承、Windows 的动态优先级提升是策略。

\Q{Q4-2.7 实例操作系统的调度算法都是什么？}
传统 UNIX：动态优先数法；4.3BSD/5.3BSD：多级反馈队列法；Windows：基于优先级的抢占式多任务调度（以线程为单位，动态优先级 + 时间配额）；Linux：抢占式调度，实时进程用 FIFO/RR，普通进程经历了 2.4 O(n) $\to$ O(1) $\to$ CFS（2.6.23）$\to$ EEVDF（6.6 以后）；Solaris：综合调度算法（TS 分时、IA 交互、SYS 系统、RT 实时、FSS 公平共享等多个调度类）；xv6：扫描进程表轮转，没有优先级；Go 运行时：GMP 模型 + 工作窃取。

\subsection{处理器调度的基本概念：层次、时机、过程}
\begin{btab}{P{0.300}P{0.483}P{0.217}}
\thead{层次 & 做什么 & 对应的状态转换}
\trow{长程调度（作业调度、宏观调度） & 创建新进程时，决定是否把它放入当前活跃进程的集合，从而控制多道程序的“道数”和系统负载 & 新建$\to$就绪}
\trow{中程调度（交换） & 从存储管理的角度，把进程的部分或全部换出到外存，为正在运行的进程腾出内存；需要时再换入 & 就绪/阻塞$\leftrightarrow$挂起}
\trow{\textbf{短程调度}（进程调度、微观调度） & 从处理机分配的角度，从就绪进程中选一个进入运行状态。时间尺度通常为毫秒级，调用非常频繁，\imp{实现必须高效} & 就绪$\to$运行}
\end{btab}
在现代通用 OS 中，长程调度基本消失了（集群的作业队列系统中仍有）；中程调度演变为以页为单位的换出，以及内存不足时杀进程（Linux 的 OOM killer、Android 的 LMK）。本章讨论的默认是短程调度。调度程序是内核中的一个函数（Linux 的 \cd{schedule()}）；xv6 是个例外，每个 CPU 有一个专门的调度器线程。

\textbf{调度时机} \bk（老师：“待会要提问，要认真记”，属于“核心 20\%”）：
\begin{enumerate}
    \item 进程执行完毕，正常退出；
    \item 进程因为某种错误或异常而终止；
    \item 新进程被创建；
    \item 正在运行的进程要等待 I/O 操作或其他资源，进入阻塞态；
    \item 被唤醒的阻塞进程重新回到就绪态（如 I/O 完成中断）；
    \item 进程用完了分配给它的时间片，回到就绪队列；
    \item 进程主动放弃 CPU（执行 \cd{yield}）。
\end{enumerate}
对应的典型事件：创建、唤醒、退出等进程控制操作；进程等待 I/O、I/O 中断；时钟中断（时间片用完、计时器到时）；进程执行中出现 abort 异常；执行 \cd{yield}。\textbf{本质}：\textbf{就绪队列发生了变化}，或者 CPU 空闲了。

\textbf{统一的流程}（课件原文）：事件发生 $\to$ 当前运行的进程暂停运行 $\to$ 硬件机制响应后进入 OS，处理相应的事件 $\to$ 处理结束后，某些进程的状态会发生变化，也可能出现了新的进程 $\to$ 就绪队列有调整 $\to$ 需要由进程调度根据预设的算法从就绪队列中选一个进程。调度发生在\textbf{内核对中断/异常/系统调用处理完、返回用户态之前的最后时刻}。

\textbf{两种视角的对应}：进程状态模型中的每一条转换边，都对应中断/异常机制中一个确定的事件。例如 \cd{fork} 创建新进程 $\leftrightarrow$ 系统调用；时间片用完 $\leftrightarrow$ 时钟中断；阻塞进程被唤醒 $\leftrightarrow$ I/O 中断。所以“进程状态转换”“中断异常处理”“调度时机”是同一件事的三个侧面。实现上，Linux 在事件处理中只置 \cd{TIF_NEED_RESCHED} 标志，在统一出口检查后调用 \cd{schedule()}；xv6 中，\cd{yield}（时间片用完）、\cd{sleep}（阻塞）、\cd{exit}（退出）最终都调用 \cd{sched()}。

\ez \imp{时钟中断 $\neq$ 时间片用完}：时钟中断是硬件按固定频率（如每 1ms 或 10ms）发出的信号；时间片是内存中的一个计数器，每次时钟中断让它减 1，减到 0 才叫“时间片用完”。所以多数时钟中断并不会引起调度。

\textbf{调度过程——进程切换}：调度程序选出的进程可以是刚被暂停的那个进程，也可以是另一个进程。\textbf{进程切换}是指一个进程让出处理器，由另一个进程占用处理器的过程，切换时要保存原进程的各种状态，恢复新进程的各种状态。进程运行时，其硬件状态保存在 CPU 的寄存器中（PC、程序状态寄存器、栈指针、通用寄存器、其他控制寄存器）；不运行时，这些值保存在 PCB 中。把 CPU 的硬件状态从一个进程换到另一个进程的过程称为\textbf{上下文切换}，主要包括两部分工作：(1) \textbf{切换全局页目录}，加载一个新的地址空间；(2) \textbf{切换内核栈和硬件上下文}，硬件上下文包括内核执行新进程所需的全部信息，如 CPU 相关寄存器。其中重点是\textbf{更换页表基址寄存器}。

xv6 的 \cd{swtch(old, new)}：把 \cd{ra}、\cd{sp}、\cd{s0}–\cd{s11} 依次存入 \cd{*old}，再从 \cd{*new} 中取出这些寄存器，最后 \cd{ret} 跳到 \cd{new->ra}。换了 \cd{sp} 就意味着换了内核栈；用户态寄存器早已保存在 trapframe 中（两层保存，见第 2 章 xv6 题 (3)）。上游 xv6 切换用户页表（写 \cd{satp}）是在返回用户态的 \cd{usertrapret}/\cd{userret} 中完成的；实验用的 xv6-k210 每个进程还有一张内核页表，所以在 \cd{scheduler} 中 \cd{swtch} 前后要 \cd{w_satp} 并 \cd{sfence_vma}。

\subsection{调度算法的设计考虑}

\textbf{性能指标}：
\begin{itemize}
    \item \textbf{周转时间}（turnaround time）= 完成时刻 $-$ 到达（提交）时刻，即进程从提交到完成所经历的总时间，包括等待和运行；\textbf{平均周转时间}是所有进程周转时间的平均值。
    \item \textbf{响应时间}（response time）= 从提出请求到\textbf{第一次}得到回应的时间。
    \item \textbf{等待时间} = 周转时间 $-$ 运行时间，即在就绪队列中等待的总时间。
    \item \textbf{吞吐量}（throughput）= 单位时间内完成的进程数。
    \item \textbf{CPU 利用率} = CPU 做有效工作的时间比例。
    \item 非性能指标：\textbf{公平性}（每个进程得到合理的 CPU 份额，不会饥饿）、\textbf{均衡性}（CPU 和 I/O 设备都忙）、\textbf{可预测性}、强制执行优先级；特定场景还有实时性、能耗。
\end{itemize}
例：编译一个大程序，1 秒后屏幕上出现第一行输出（响应时间 1s），60 秒后全部编译完（周转时间 60s）。不同系统看重的指标不同：交互式进程要求响应时间短（50–150ms），批处理看重吞吐量和周转时间，实时系统看重截止时间和可预测性。

\textbf{PCB 中需要记录哪些与调度有关的信息？} 调度的依据\textbf{必须在 PCB 中有记录}，否则调度程序无从判断。包括：进程状态、优先级（静态/动态）、调度策略或调度类、剩余时间片或已运行时间、等待/睡眠时间的统计、允许运行在哪些 CPU 上（亲和性）、队列指针、保存的上下文。xv6 的 PCB 中与调度相关的只有 \cd{state}；Windows 线程有\textbf{基本优先级}和\textbf{当前优先级}两个字段。

\textbf{优先级与优先数} \ez：\textbf{优先级}表示进程的重要性和紧迫程度，是一个概念；\textbf{优先数}是用来量化优先级的数值。答题时\imp{要先说清方向}：UNIX/Linux 的 nice 值为 $-20$（最高）到 19（最低），内核内部的 \cd{prio} 为 0–139，数值越小优先级越高；Windows 的优先级为 0–31，数值越大优先级越高。

\textbf{静态优先级与动态优先级}：静态优先级在进程创建时指定，运行过程中不再改变；动态优先级在创建时指定一个初值，运行过程中可以动态变化。动态调整的常见规则：等待时间越长，优先级越高（称为\textbf{老化} aging，防止饥饿）；使用 CPU 越多，优先级越低（照顾 I/O 密集型进程）。

\textbf{就绪队列的组织}：按优先级设置多个就绪队列，调度时总是从优先级最高的非空队列中取进程。为了快速找到“最高的非空队列”，可以用一个\textbf{位图}记录哪些队列非空，再用“查找第一个置位位”的指令在 O(1) 时间内找到（Windows 的就绪位图、Linux O(1) 调度器都这样做）。还可以允许进程在队列之间移动（这就是多级反馈队列的雏形）。

\textbf{I/O 密集型与 CPU 密集型进程}：为什么调度要\textbf{偏好 I/O 密集型进程}？\bk 两条理由：
\begin{enumerate}
    \item I/O 密集型进程只用一小段 CPU 就会发出 I/O 请求、让出 CPU。优先调度它，能让它尽快把 I/O 发出去，使 \textbf{CPU 和外设同时忙碌}，提高整个系统的资源利用率。
    \item 它的时间片往往用不完就阻塞了，优先调度它可以看作是一种补偿，体现公平。
\end{enumerate}
随着 CPU 越来越快，而磁盘、网络相对越来越慢，进程越来越倾向于 I/O 密集型，这个问题也就越来越重要。

\textbf{时间片大小的选择}：需要考虑进程切换的开销、对响应时间的要求、就绪进程的个数、CPU 的处理能力、进程的行为。
\begin{itemize}
    \item \imp{时间片太长}：每个进程都能在一个时间片内运行完或阻塞，RR \textbf{退化为 FCFS}，交互进程的响应时间变长。
    \item \imp{时间片太短}：切换次数多，切换开销所占比例升高；而且一次交互要跨越多个时间片才能完成，响应时间反而变长。
    \item 经验：时间片应略大于一次典型交互所需的 CPU 时间。若时间片为 10ms、一次切换耗时 0.1ms，开销约为 1\%；若时间片为 1ms，开销约为 9\%。
    \item 例：一个需要运行 10 个时间单位的进程，时间片为 12 时切换 0 次，为 6 时切换 1 次，为 1 时切换 9 次。
    \item 实际系统：Linux \cd{SCHED_RR} 默认 100ms；Windows 客户端版约为 2 个时钟间隔；上游 xv6 在每次时钟中断时都 \cd{yield}，即时间片等于一个时钟间隔。
\end{itemize}

\subsection{批处理系统中的调度算法}

\textbf{先来先服务 FCFS}（First Come First Serve，又称 FIFO）：按照进程进入就绪队列的先后顺序选择，是\textbf{非抢占式}的：选中的进程一直运行到结束或阻塞。阻塞的进程被唤醒后排到队尾。优点：公平（按先后），实现简单。缺点：短进程排在长进程后面要等很久（\textbf{护航效应}，convoy effect）；对 I/O 密集型进程不利，它们每次只用一点点 CPU 却要等整个队列，I/O 设备经常闲置。

例（课件）：P1、P2、P3 依次就绪，P1 需要 24s，P2、P3 各需要 3s。按 P1、P2、P3 的顺序：周转时间分别为 24、27、30，\imp{平均周转时间 27s}；如果改变顺序为 P2、P3、P1：周转时间为 3、6、30，平均只有 13s。两种顺序的吞吐量都是 3 jobs / 30s = 0.1 jobs/s。可见短进程先运行能大大降低平均周转时间——这就引出了 SJF。

\textbf{短作业优先 SJF}（Shortest Job First）：选择运行时间最短的进程先运行，是\textbf{非抢占式}的：选中的进程一直运行到结束或阻塞。\textbf{最短剩余时间优先 SRTN}（Shortest Remaining Time Next）是它的\textbf{抢占式}版本：当一个新进程就绪时，如果它需要的运行时间比当前进程的\textbf{剩余}运行时间还短，就抢占当前进程。

例：P1（到达时刻 0，运行时间 7）、P2（2，4）、P3（4，1）、P4（5，4）。
\begin{itemize}
    \item \textbf{SJF}：t=0 只有 P1，运行到 7；t=7 时 P2、P3、P4 都已到达，选最短的 P3（1），t=8 完成；P2 和 P4 一样长，按先来选 P2，t=12 完成；P4 在 t=16 完成。执行序列 P1[0–7] P3[7–8] P2[8–12] P4[12–16]。周转时间 7、10、4、11，\textbf{平均 8}；等待时间 0、6、3、7，平均 4。
    \item \textbf{SRTN}：t=2 时 P2 到达，它需要 4，而 P1 剩余 5，于是 P2 抢占 P1；t=4 时 P3 到达，需要 1，P2 剩余 2，P3 抢占；t=5 P3 完成，此时 P4 到达需要 4，P2 剩余 2，P2 继续；t=7 P2 完成，P4 运行到 11；最后 P1 运行剩下的 5，t=16 完成。执行序列 P1[0–2] P2[2–4] P3[4–5] P2[5–7] P4[7–11] P1[11–16]。周转时间 16、5、1、6，\textbf{平均 7}；等待时间 9、1、0、2，平均 3。
\end{itemize}
\textbf{优点}：\textbf{平均周转时间最短}（所有作业同时到达时，SJF 是最优的）。证明用\textbf{交换论证}：假设相邻两个作业 $a$ 在 $b$ 之前，且 $s_a > s_b$；交换它们后，$b$ 的完成时间提前 $s_a$，$a$ 的完成时间推后 $s_b$，其他作业不受影响，总周转时间减少 $s_a - s_b > 0$。所以只要存在“长在短前”，就能通过交换变得更好，最优顺序必然是从短到长。

\textbf{缺点}：(1) 长作业可能\textbf{饥饿}，只要不断有短作业到来，长作业就一直得不到运行，不公平；(2) 需要\textbf{预知运行时间}，这在实际中做不到。批处理系统可以让用户申报估计的运行时间；交互式系统则用历史数据预测下一次 CPU 突发的长度，常用\textbf{指数平均}：$\tau_{n+1} = \alpha t_n + (1-\alpha)\tau_n$，其中 $t_n$ 是第 $n$ 次实际运行时间，$\tau_n$ 是对它的预测。例：$\alpha = 1/2$，$\tau_0 = 10$，实际运行时间依次为 6、4、6、4，则预测值依次为 8、6、6、5。这就是金融中的 EWMA（指数加权移动平均）。用预测值调度的版本就是交互式系统中的\textbf{最短进程优先 SPN}（见 4.5 节）。

\textbf{最高响应比优先 HRRN}（Highest Response Ratio Next）：每次调度时计算所有就绪进程的响应比，选最高的。
\[ R = \frac{\text{周转时间}}{\text{处理时间}} = \frac{\text{处理时间} + \text{等待时间}}{\text{处理时间}} = 1 + \frac{\text{等待时间}}{\text{处理时间}} \]
课件的思考：
\begin{itemize}
    \item \textbf{短作业}：处理时间短，在等待时间相同的情况下响应比更大，更容易被调度——类似 SJF，降低平均周转时间。
    \item \textbf{长作业}：等待时间足够长后，它的响应比会超过后来的短作业，从而得到运行机会——\textbf{避免饥饿}。
    \item 处理时间相同时，等得越久响应比越高——类似 FCFS。
\end{itemize}
所以 HRRN 是 FCFS 与 SJF 的\textbf{折中}。例：长作业 $s=20$ 已等 60，$R = 1 + 60/20 = 4$；短作业 $s=1$ 已等 2，$R = 1 + 2/1 = 3$；此时选长作业。HRRN 通常是非抢占式的。如果做成抢占式，需要频繁地重新计算所有进程的响应比并排序（等待时间时刻在变），开销太大，所以实际中很少使用。

\subsection{交互式系统中的调度算法}

\textbf{时间片轮转 RR}（Round Robin）：为每个进程分配一个时间片；当时间片用完时，时钟中断发生，OS 把当前进程放到就绪队列\textbf{末尾}，选队首的进程运行；进程在时间片内阻塞或结束，则立即调度。它是抢占式的（按时间片抢占）。

例：时间片 $q=20$，P1=53、P2=8、P3=68、P4=24，同时到达。执行序列：P1[0–20] P2[20–28] P3[28–48] P4[48–68] P1[68–88] P3[88–108] P4[108–112] P1[112–125] P3[125–145] P3[145–153]。等待时间：P1 $= (68-20)+(112-88) = 72$，P2 $= 20$，P3 $= 28+40+17 = 85$，P4 $= 48+40 = 88$，\textbf{平均 66.25}。

同一组数据在不同时间片下的平均等待时间：$q=1$ 为 62，$q=5$ 为 61.25，$q=8$ 为 57.25，$q=10$ 为 61.25，$q=20$ 为 66.25；最好的 FCFS 顺序（P2、P4、P1、P3，即从短到长）为 31.25，最坏的 FCFS 顺序为 83.5。观察：RR 的结果介于最好和最坏的 FCFS 之间；P3 最长，无论 $q$ 取多少都是最后完成；$q$ 恰好等于 P2 的长度 8 时结果最好。

\textbf{优缺点}：公平，响应时间短，适合交互式和分时系统；缺点是切换带来额外开销，并且\imp{对长度相同的进程不利}。例：A、B 两个进程各需 100ms，时间片 1ms。RR 下两者交替运行，分别在 199ms 和 200ms 完成，平均周转时间 199.5ms；FCFS 下分别在 100ms 和 200ms 完成，平均 150ms。RR 用周转时间换取了响应时间。

\textbf{课堂练习：RR 的 CPU 利用率} \bk。设每个进程运行 $T$ 时间后就阻塞去做 I/O，一次进程切换的开销为 $S$，时间片为 $Q$，并假设总有就绪进程。思路：CPU 利用率 $=$ 有用时间 $/$（有用时间 $+$ 切换时间），关键是看\textbf{运行完一段 $T$ 需要切换几次}（进程阻塞时也要切换一次）。
\begin{enumerate}
    \item $Q = \infty$：退化为 FCFS，每段 $T$ 运行完阻塞时切换 1 次，利用率 $\dfrac{T}{T+S}$。
    \item $Q > T$：时间片还没用完进程就阻塞了，同样每段切换 1 次，利用率 $\dfrac{T}{T+S}$。
    \item $S < Q < T$：每段 $T$ 被切成多个时间片，每个时间片后跟一次切换，利用率约为 $\dfrac{Q}{Q+S}$。精确值为 $\dfrac{T}{T+S\lceil T/Q\rceil}$，例如 $T=10$、$Q=4$、$S=1$ 时，精确值为 $10/13 \approx 76.9\%$，近似值为 80\%。
    \item $Q = S$：代入上式，利用率为 \imp{50\%}。
    \item $Q \to 0$：几乎所有时间都在切换，利用率趋于 0。
\end{enumerate}
结论：利用率随 $Q$ 增大而提高，$Q \ge T$ 以后封顶为 $\dfrac{T}{T+S}$。$Q$ 太小开销大，太大则退化为 FCFS、响应差，所以 $Q$ 应略大于典型的 CPU 突发长度，并远大于 $S$。

\textbf{虚拟轮转 VRR}：RR 对 I/O 密集型进程不公平——它们每次只用了时间片的一小部分就去等 I/O，回来后却要排到队尾，而 CPU 密集型进程每次都用满整个时间片。VRR 的改进：设置一个\textbf{辅助队列}，I/O 完成的进程进入辅助队列，\textbf{辅助队列优先于}普通就绪队列被调度；从辅助队列选出的进程只运行\textbf{它上次时间片剩余的部分}。这样 I/O 密集型进程得到补偿，而 CPU 密集型进程照常轮转，不受影响。

\textbf{混合负载例（课堂派）}：A、B 是 CPU 密集型进程，各需运行 1000ms；C 是 I/O 密集型进程，循环执行“1ms CPU + 10ms 磁盘 I/O”（如果单独运行，磁盘利用率为 $10/11 \approx 91\%$）。目标是让 CPU 和磁盘都尽量忙。
\begin{itemize}
    \item \textbf{FCFS}：如果 A、B 排在 C 前面，C 要等 2000ms 才能发出第一次 I/O，这期间\textbf{磁盘一直闲置}。
    \item \textbf{RR，$q=100$ms}：C 每运行 1ms 就去做 I/O，回来后要等 A、B 各用完 100ms，所以每 200ms 磁盘只工作约 10ms，利用率约 5\%。
    \item \textbf{RR，$q=1$ms}：C 的 I/O 完成后很快就能被调度，磁盘利用率\textbf{接近 90\%}，A、B 几乎不受影响；代价是\textbf{切换次数极多}。
    \item \textbf{SRTN}：C 的下一次 CPU 突发只有 1ms，总是最短的，所以 I/O 一完成就抢占 A 或 B，磁盘利用率高，切换次数也少；但需要预测运行时间，而且 B 可能饥饿。
\end{itemize}
结论：照顾 I/O 密集型进程（CPU 突发短的进程），能同时提高 CPU 和 I/O 设备的利用率。

\textbf{最高优先级 HPF}（Highest Priority First）：选择优先级最高的进程运行。一般原则：系统进程高于用户进程，前台进程高于后台进程，偏好 I/O 密集型进程。优先级可以是静态的或动态的，调度可以是抢占式或非抢占式的。实现简单，但\textbf{不公平，低优先级进程可能饥饿}（传说 MIT 的 IBM 7094 在 1973 年关机时，发现一个 1967 年提交的低优先级作业还没有运行）。解决办法是老化：等待越久，优先级越高。

\textbf{优先级反转} \bk：
\begin{itemize}
    \item \textbf{现象}：低优先级进程 L 持有高优先级进程 H 所需要的资源（如一把锁），H 因此阻塞；这时一个中优先级进程 M（不需要这个资源）就绪，抢占了 L 并一直运行。L 得不到 CPU，就无法释放锁，H 只能一直等下去。结果是 H 的等待时间\textbf{没有上界}，M 实际上“压过”了 H。
    \item \textbf{前提}：\imp{只出现在基于优先级的抢占式调度中}。在非抢占式调度下，L 不会被 M 抢占，它会很快用完锁并释放。
    \item \textbf{影响}：系统错误、高优先级进程停滞不前、系统性能下降。
    \item \textbf{案例：火星探路者}（1997 年，VxWorks 实时系统）：高优先级的总线管理任务、中优先级的通信任务、低优先级的气象数据采集任务。气象任务持有总线的互斥锁时被通信任务抢占，总线管理任务因拿不到锁而超时，看门狗判定系统故障，导致探测器\textbf{反复重启}。VxWorks 的互斥锁其实有“优先级继承”选项，但默认是关闭的；工程师在地面复现问题后，远程打开了该选项，问题得以解决。
    \item \textbf{解决方法}：(1) \textbf{优先级继承}：当 H 因为 L 持有的锁而阻塞时，L 临时继承 H 的优先级，这样 M 就不能抢占 L；L 释放锁后恢复原优先级。(2) \textbf{优先级天花板}：为每把锁设定一个“天花板”优先级，等于可能使用它的进程中最高的优先级；任何进程拿到这把锁，就把优先级提升到天花板。(3) 进入临界区时\textbf{关中断}（只适用于单 CPU 且临界区很短的情况）。Linux 提供 \cd{PTHREAD_PRIO_INHERIT}、\cd{PTHREAD_PRIO_PROTECT} 两种互斥锁协议，内核中有 rt\_mutex。
\end{itemize}

\textbf{多级队列}：根据进程的类型把它们\textbf{永久地}分到不同的就绪队列中（如系统进程、交互进程、编辑进程、批处理进程、学生进程），进程不会在队列之间移动。队列之间按固定优先级调度，或按比例分配 CPU 时间（如前台 80\%、后台 20\%）；每个队列可以采用不同的调度算法（前台用 RR，后台用 FCFS）。Linux 的调度类 deadline $>$ rt $>$ fair $>$ idle 实际上就是多级队列。

\textbf{多级反馈队列 MLFQ} \bk（Corbató 1962 年在 CTSS 中提出，他因此获得图灵奖；UNIX BSD 5.3 采用）：它\textbf{不需要预先知道}进程的类型和运行时间，而是根据进程的\textbf{历史行为（反馈）}动态调整它所在的队列。规则：
\begin{enumerate}
    \item 设置多个就绪队列，第一级队列的优先级最高。
    \item 各级队列的时间片不同，\imp{优先级越低，时间片越长}（如 8ms、16ms、32ms……）。
    \item 只有高级队列为空时，才调度下一级队列中的进程。
    \item 同一级队列内按 RR 调度（最后一级可以用 FCFS）。
    \item \textbf{新创建的进程进入第一级队列}。
    \item \textbf{进程用完时间片后放弃 CPU，降到下一级队列}（已在最底层则不再降）。
    \item \textbf{进程因阻塞而放弃 CPU，等待的事件发生后回到原来那一级队列，不降级}——这一条使 MLFQ 偏好 I/O 密集型进程。
    \item 在允许抢占的情况下，被更高优先级进程抢占的进程回到原来那一级队列（通常放在队首）。
\end{enumerate}
\textbf{课件中打问号的两处就是考点}（老师说“没有唯一标准答案，属于策略”，答题时要写出各种选项和理由）：
\begin{itemize}
    \item \textbf{解除阻塞、被抢占后，放回原队列的队首还是队尾？} 被抢占的进程放\textbf{队首}更公平，因为它的时间片是被别人打断的，并非自己用完（Windows 就是这样做的）；放队尾实现更简单。阻塞后被唤醒的进程一般放队尾；也可以放回原级、升一级（更照顾 I/O 密集型），或者直接回到第一级（响应最快，但可能被恶意利用）。
    \item \textbf{再次被调度时，使用剩余的时间片、新的完整时间片，还是奖励时间片？} 被抢占的进程可以只用剩余的时间片（Windows 中可变优先级线程就是这样），也可以给一个新的完整时间片；阻塞后被唤醒的进程通常给一个完整的时间片；给奖励时间片更照顾 I/O 密集型进程，但可能被钻空子。需要在公平性、I/O 偏好和开销之间权衡。
\end{itemize}
\textbf{MLFQ 的问题和修正}（OSTEP）：(1) \textbf{饥饿}：如果交互进程很多，它们一直占据高优先级队列，CPU 密集型进程沉在底层永远得不到运行。修正：每隔一段时间做一次\textbf{优先级提升}（boost），把所有进程移回最高级队列；这也解决了进程行为发生变化（如从计算阶段转入交互阶段）的问题。(2) \textbf{愚弄调度器}：恶意进程在时间片用完前的一刻发出一个无意义的 I/O，就能永远不降级。修正：\textbf{累计记账}，记录进程在本级队列中用掉的总时间，用满本级的额度就降级，无论中间是否放弃过 CPU。(3) 参数（队列数、各级时间片、提升周期）很难调好。MLFQ 的实质是：\textbf{在不知道运行时间的情况下近似实现 SJF/SRTN}——短进程在高级队列中很快就完成了，长进程逐渐沉到下面。

\textbf{最短进程优先 SPN}（Shortest Process Next）：把 SJF 的思想搬到交互式系统中。交互式进程的行为是“等待命令 $\to$ 执行一小段 CPU（一次 CPU 突发）$\to$ 再等待”，把\textbf{每一次 CPU 突发看作一个独立的“作业”}，每次调度时选择\textbf{预计下一次 CPU 突发最短}的就绪进程。
\begin{itemize}
    \item \textbf{难点：怎么知道下一次突发有多长？} 用该进程过去的突发长度来\textbf{预测}，常用\textbf{指数平均}（老化）：$\tau_{n+1} = \alpha t_n + (1-\alpha)\tau_n$，$0 \le \alpha \le 1$。$\alpha$ 越大越看重最近一次，$\alpha = 0$ 则永远用初值，$\alpha = 1$ 则只看上一次。展开后 $\tau_{n+1} = \alpha t_n + (1-\alpha)\alpha t_{n-1} + (1-\alpha)^2 \alpha t_{n-2} + \cdots$，越早的历史权重按几何级数衰减（即金融中的 EWMA）。
    \item \textbf{例}：$\alpha = 1/2$，初始预测 $\tau_0 = 10$，实际突发依次为 6、4、6、4，预测值依次为 $\tau_1 = (6+10)/2 = 8$，$\tau_2 = (4+8)/2 = 6$，$\tau_3 = (6+6)/2 = 6$，$\tau_4 = (4+6)/2 = 5$。预测值逐渐向实际值靠拢。
    \item \textbf{为什么适合交互式系统}：交互进程（编辑器、shell）每次只用一点点 CPU 就去等用户输入，预测值小，会被优先调度，\textbf{响应时间短}；这也就是“偏好 I/O 密集型进程”。
    \item \textbf{缺点}：(1) 和 SJF 一样，CPU 突发长的进程可能\textbf{饥饿}；(2) 预测可能不准，进程行为突然改变时要过几轮才能跟上；(3) 要为每个进程记录历史并计算预测值，有额外开销；(4) 可以做成非抢占的，也可以做成抢占的（新就绪进程的预测值小于当前进程的剩余预测值时抢占，即 SRTN 的思想）。
    \item \textbf{与 MLFQ 的关系}：SPN 是“显式地预测突发长度”，MLFQ 是“根据是否用完时间片隐式地推断突发长度”——二者都是在不知道运行时间时近似 SJF。
\end{itemize}

\textbf{按比例分配的调度算法}：\par
\begin{btab}{P{0.109}P{0.446}P{0.446}}
\thead{算法 & 原理 & 评价}
\trow{公平共享 & 公平针对的是\textbf{用户或用户组}而不是进程：先在用户之间平分 CPU，再在每个用户的进程之间分配。例：用户 1 有进程 A、B、C、D，用户 2 只有进程 E，调度序列为 A E B E C E D E & 防止某个用户开很多进程抢占资源；需要维护按用户的统计。Linux cgroup 的 \cd{cpu.weight}、Solaris 的 FSS 调度类}
\trow{保证调度 & 向用户做出明确的性能保证：有 n 个进程时，每个进程获得 1/n 的 CPU 时间。系统跟踪每个进程“应得”与“实得”的 CPU 时间，选比值最低的运行 & 需要持续记账；CFS 的 vruntime 就是这一思想}
\trow{彩票调度 & 给进程发彩票，每次调度时随机抽一张，持有该彩票的进程获得 CPU；进程持有彩票的比例就是它期望获得的 CPU 份额。进程可以\textbf{转让}彩票（如客户端阻塞等待服务器时，把彩票借给服务器），也可以“通货膨胀”（自己增发彩票） & 简单，能按比例分配，新进程拿到彩票立即参与；但基于概率，短期内可能不公平。它的确定性版本是\textbf{步长调度}：步长 = 一个大数 / 票数，每次选 pass 值最小的进程运行，运行后 pass 加上步长——CFS 的 vruntime 就相当于 pass}
\end{btab}

\textbf{线程调度}：
\begin{itemize}
    \item \textbf{用户级线程}：内核只调度进程，线程库在进程内部调度线程，时间片是分给进程的。例：进程 A 有线程 A1、A2、A3，进程 B 有 B1、B2、B3，进程时间片 50ms，每个线程每次运行 5ms。可能的序列是 A1 A2 A3 A1 A2 A3……（在 A 的 50ms 内），然后才轮到 B 的线程；\imp{不可能出现 A1 B1 这样的交替}，因为内核在 A 的时间片内不知道有 B1。优点是切换快，调度策略可以由应用定制；缺点是一个线程阻塞则整个进程阻塞。
    \item \textbf{核心级线程}：内核直接调度线程，可能出现 A1 B1 A2 B2 A3 B3 这样的序列，也可能是 A1 A2 A3（取决于调度器）。在同一进程的线程之间切换不需要换地址空间，比切换到另一个进程的线程更便宜，内核调度时可以考虑这一因素。
\end{itemize}

\subsection{调度算法的比较}

\textbf{综合对比}（$w$ = 已花费的等待时间，$e$ = 到现在为止已花费的执行时间，$s$ = 进程所需的总服务时间，包括 $e$）：\par
\begin{btab}{P{0.109}P{0.130}P{0.232}P{0.232}P{0.232}P{0.065}}
\thead{算法 & 选择函数 & 决策模式 & 吞吐量 / 响应时间 & 开销与影响 & 饥饿}
\trow{FCFS & $\max[w]$ & 非抢占 & 不强调 / 响应时间可能很长，特别是进程执行时间差别大时 & 开销最小；对短进程、I/O 密集型进程不利 & 无}
\trow{RR & 常数 & 抢占（时间片用完时） & 时间片太小时吞吐量低 / 短进程响应时间好 & 开销最小；公平对待所有进程 & 无}
\trow{SJF & $\min[s]$ & 非抢占 & 高 / 短进程响应好 & 开销可能较高；对长进程不利 & 可能}
\trow{SRTN & $\min[s-e]$ & 抢占（新进程到达时） & 高 / 好 & 开销可能较高；对长进程不利 & 可能}
\trow{HRRN & $\max[(w+s)/s]$ & 非抢占 & 高 / 好 & 开销可能较高；很好的平衡 & 无}
\trow{MLFQ & 见上文规则 & 抢占（时间片用完时） & 不强调 / 不强调 & 开销可能较高；可能对 I/O 密集型进程有利 & 可能}
\end{btab}
记忆规律：选择函数中含有等待时间 $w$ 的算法不会饿死进程（等得越久越容易被选中）；只看 $s$ 的算法会饿死长进程；RR 靠轮转保证不饿死；MLFQ 可能饿死底层的 CPU 密集型进程。

\subsection{实例：Windows 线程调度}

\textbf{调度单位是线程}，采用\textbf{基于动态优先级的、抢占式}调度，结合\textbf{时间配额}的调整。具体做法：就绪线程按优先级进入相应的队列；系统总是选择优先级最高的就绪线程运行；同一优先级的各线程按时间片轮转；多处理机系统中允许多个线程并行运行。

\textbf{引发调度的条件}：线程创建或变为就绪、线程终止、线程进入等待（阻塞）、线程从运行变为就绪（时间配额用完或被抢占），以及 Windows 特有的两个：\textbf{线程的优先级改变}、\textbf{线程的亲和处理机集合改变}（允许运行的 CPU 变了，当前 CPU 可能不再允许它运行）。

\textbf{32 个优先级}（数值越大越高）：
\begin{itemize}
    \item 16–31：\textbf{实时优先级}，线程的优先级不会被系统改变。
    \item 1–15：\textbf{可变优先级}。线程的\textbf{基本优先级}由进程的优先级类别（IDLE 为 4、NORMAL 为 8、HIGH 为 13、REALTIME 为 24 等）加上线程的相对优先级决定；\textbf{当前优先级}可以在基本优先级之上临时提升，之后逐步衰减回基本优先级。
    \item 0：系统的零页线程（把空闲页清零）。
\end{itemize}

\textbf{时间配额}：不是一个时间长度，而是一个\textbf{配额单位}的整数。每次时钟中断，当前线程的配额减 3（\cd{CLOCK_QUANTUM_DECREMENT}）；线程完成一次等待也减 1。客户端版本的 \cd{QuantumReset} 默认为 6（即 2 个时钟间隔），服务器版本为 36（12 个时钟间隔，减少切换、提高吞吐）。配额用完时，如果没有其他同优先级的就绪线程，就重新给它一份配额继续运行。

\textbf{为什么给前台进程增加配额，而不是提高优先级？} \textbf{优先级决定谁先运行，配额决定一次运行多久}。如果把前台游戏的优先级提高，后台的电子表格计算可能完全得不到 CPU（饿死）；而增加游戏的配额，表格计算只是慢一些，仍能运行。

\textbf{调度数据结构}：每个优先级一个就绪队列；一个 32 位的\textbf{就绪位图}，每一位表示对应优先级的队列是否非空，可以在 O(1) 时间内找到最高优先级；一个\textbf{空闲位图}，每个 CPU 一位，表示该 CPU 是否空闲。

\textbf{调度策略}：
\begin{itemize}
    \item \textbf{主动切换}：线程进入等待状态（如等待 I/O），调度程序选择优先级最高的就绪线程运行。
    \item \textbf{抢占}：一个更高优先级的线程变为就绪，当前线程被抢占，\textbf{回到其优先级队列的队首}（而不是队尾）。如果被抢占的是\textbf{实时}线程，它的配额被\textbf{重置为完整的一份}；如果是\textbf{可变优先级}线程，它的\textbf{配额保持不变}，下次运行时继续用剩余的配额。
    \item \textbf{配额用完}：如果线程的优先级没有降低，且有同优先级的其他就绪线程，就把它放到队尾，选下一个线程；如果没有同优先级的就绪线程，就重新分配配额让它继续运行。如果线程的优先级降低了（之前被提升过，现在衰减），就选择更高优先级的线程运行。\ez \imp{配额用完并不会降低优先级}——这一点与 MLFQ 不同；只有被提升过的线程才会逐级衰减回基本优先级。
\end{itemize}

\textbf{五种优先级提升}（只针对 1–15 的可变优先级线程）：
\begin{enumerate}
    \item \textbf{I/O 操作完成}：提升幅度由设备驱动程序在 \cd{IoCompleteRequest} 中指定（键盘、声卡提升多，磁盘提升少），同时配额减 1。
    \item \textbf{信号量或事件等待结束}：提升 1 级（不超过 15），配额减 1；之后每用完一个配额就降 1 级，直到回到基本优先级。
    \item \textbf{前台进程中的线程完成一次等待}：提升，以改善交互响应。
    \item \textbf{由于窗口活动而唤醒 GUI 线程}：提升。
    \item \textbf{线程处于就绪状态太久而没有运行（饥饿）}：\textbf{平衡集管理器}每秒扫描一次就绪队列，发现某线程已就绪超过 \textbf{300 个时钟中断间隔}（约 4 秒）仍未运行，就把它的优先级提升到 \textbf{15}，并给它 \textbf{4 倍的配额}；用完配额后立即降回原来的优先级。这一机制顺带也缓解了优先级反转问题。
\end{enumerate}
\textbf{空闲线程}：每个 CPU 有一个，优先级为 0。它循环执行：处理挂起的中断 $\to$ 处理 DPC（延迟过程调用，相当于 Linux 的 softirq，即中断下半部）$\to$ 检查是否有就绪线程，有则切换过去 $\to$ 调用 HAL 中的处理器空闲例程（省电）。

\subsection{多处理器调度}

多处理器上，调度程序不仅要决定选哪个进程，还要决定在哪个 CPU 上运行；需要考虑进程在多个 CPU 之间迁移时的开销（cache、TLB 失效），以及如何使所有 CPU 的负载均衡。多处理器分为 \textbf{SMP}（对称多处理，每个 CPU 地位相同，都运行自己的调度程序，访问共享的资源时需要同步）和 \textbf{NUMA}（非一致内存访问，访问本地内存比远端内存快）。\par
\begin{btab}{P{0.109}P{0.446}P{0.446}}
\thead{& 静态分配（每 CPU 一个队列，MQMS） & 动态分配（全局共享队列，SQMS）}
\trow{做法 & 进程从开始到结束都固定在一个处理器上，每个处理器有自己的就绪队列 & 所有就绪进程放在一个公共队列中，可以被分配到任何空闲的处理器上执行}
\trow{调度开销 & 小，各 CPU 调度时不需要抢同一把锁 & 大，每次调度都要锁住公共队列，CPU 越多越容易成为瓶颈}
\trow{负载 & 可能忙闲不均，需要额外的负载均衡 & 天然均衡}
\trow{亲和性 & 好，进程一直在同一个 CPU 上，cache 保持热 & 差，进程可能在 CPU 之间跳来跳去}
\end{btab}
\textbf{现代 Linux/openEuler 的做法}：每个 CPU 一个运行队列，再加上周期性的负载均衡。迁移进程依靠每个 CPU 上的\textbf{迁移线程} \cd{migration/N}（属于最高的 stop 调度类，不会被抢占）。例如 CPU0 要从 CPU1 拉一个进程过来：它向 CPU1 的停机工作队列中加入一个迁移函数，并唤醒 CPU1 的迁移线程，由 CPU1 自己把进程挪到 CPU0 的队列。之所以要让 CPU1 来做，是因为如果该进程正在 CPU1 上运行，它的上下文在 CPU1 的寄存器里，别的 CPU 拿不到。课件中还有一个“单队列 + 亲和度”的例子：进程 A–D 分别固定在各自的 CPU 上保持亲和性，只让进程 E 在 CPU 之间轮流迁移以平衡负载，即牺牲 E 的亲和性。

\textbf{多核调度要考虑的三个问题}：
\begin{itemize}
    \item \textbf{缓存一致性}：每个核有自己的 cache，同一数据可能在多个 cache 中有副本。硬件用 \textbf{MESI} 协议（Modified、Exclusive、Shared、Invalid 四种状态）保证一致：一个核要写处于 Shared 状态的缓存行时，先广播作废其他核中的副本。注意 MESI 只保证单个缓存行的一致性，并\textbf{不能}让 \cd{X = X + 1} 这样“读—改—写”的操作变成原子的，仍然需要原子指令或锁。另一个问题是\textbf{伪共享}：两个核分别频繁修改两个不同的变量，但它们恰好在同一个 64 字节的缓存行中，这一行就会在两个核之间来回失效（乒乓），性能大幅下降；解决办法是填充对齐，让它们位于不同的缓存行。
    \item \textbf{缓存亲和性}：进程在某个 CPU 上运行过一段时间后，它的数据已经在这个 CPU 的 cache 中，再次调度时最好仍放在这个 CPU 上。软亲和：调度器尽量这样做；硬亲和：用户用 \cd{taskset} 或 \cd{sched_setaffinity} 把进程绑定到指定的 CPU（例如高频交易系统把关键线程绑核，并用 \cd{isolcpus} 让这些核不运行其他任务）。
    \item \textbf{负载均衡}：避免一些 CPU 很忙、另一些空闲。
\end{itemize}
\textbf{课件思考：缓存亲和性调度与负载均衡的矛盾}。\imp{两者天然矛盾}：负载均衡要把进程迁移到空闲的 CPU 上，让所有 CPU 都干活；而迁移后，进程在原 CPU 上积累的 cache 内容就浪费了，需要在新 CPU 上重新“预热”。所以只有当迁移带来的收益大于 cache 失效的代价时才应该迁移。Linux 用\textbf{调度域}来处理：优先在“近”的 CPU 之间迁移（同一物理核的超线程之间 $\to$ 共享 L3 的核之间 $\to$ 同一 NUMA 节点内 $\to$ 跨节点），越远的迁移门槛越高。生活中的类比：并行渲染（把数据按区域分给各核以利用局部性 vs 某些区域特别复杂时要把任务转移给空闲核）；CDN（就近服务命中缓存 vs 节点过载时转发到其他节点，导致回源）。

\textbf{工作窃取（work stealing）}：任务分配到各个处理器的本地队列中，处理器优先执行本地队列中的任务；本地队列空了，就去其他处理器的队列中\textbf{偷}任务（通常偷一半）。本地队列常用双端队列实现：拥有者从一端以 LIFO 方式取任务（cache 友好），窃取者从另一端以 FIFO 方式偷（减少冲突）。
\begin{itemize}
    \item \textbf{优点}（课件思考题）：去中心化，没有全局锁；只在需要时才均衡负载；大多数时候执行本地任务，局部性好。
    \item \textbf{缺点}：偷取时需要同步，有开销；被偷走的任务失去了亲和性；偷谁的、偷多少需要策略；偷不到时会浪费时间；任务粒度太小或任务之间有依赖时效果差。
    \item \textbf{应用}：\textbf{Go 的 GMP 调度器}——G 是 goroutine，M 是内核线程，P 是逻辑处理器（数量等于 GOMAXPROCS）；M 必须持有一个 P 才能运行 G；每个 P 有一个无锁的本地 G 队列，空了就从其他 P 那里偷一半；当一个 G 执行阻塞的系统调用时，它所在的 M 被阻塞，P 会被交给另一个 M 继续运行其他 G，这就解决了 N:1 模型“一个阻塞全部阻塞”的问题。这是一种两级调度（Go 运行时调度 G 到 M 上，OS 调度 M 到 CPU 上）。此外还有 Java 的 ForkJoinPool、Intel TBB、Rust 的 Tokio。集群层面的 YARN 也有类似的考虑：数据本地性（相当于亲和性）+ ResourceManager 的负载均衡 + 多队列（Capacity/Fair 调度器）。
\end{itemize}

\subsection{实时调度}

\textbf{实时系统}要求在规定的时间内对外部事件作出反应。\textbf{硬实时}：必须满足截止时间，否则后果严重；\textbf{软实时}：偶尔错过截止时间是可以接受的。实时任务按到达方式分为\textbf{周期性}任务（每隔固定时间到达一次）、偶发任务和非周期任务。实时系统追求的是确定性，关注的是\textbf{最坏情况}下能否满足截止时间，而不是平均性能。

\textbf{可调度性}：设有 $m$ 个周期任务，任务 $i$ 的周期为 $P_i$、每次需要的 CPU 时间为 $C_i$，则能够调度的\textbf{必要条件}是总利用率不超过 1：$\sum_{i=1}^{m} C_i/P_i \le 1$。例：三个周期任务的周期分别为 100、200、500ms，每次处理时间分别为 50、30、100ms，利用率为 $0.5 + 0.15 + 0.2 = 0.85 \le 1$，满足条件；如果再加一个周期为 1000ms 的任务，它每次的处理时间最多为 150ms。

\textbf{速率单调调度 RM}（Rate Monotonic，Liu \& Layland 1973）：\textbf{静态优先级}算法。前提：任务是周期性的；每个任务的执行时间已知且不变；截止时间就是周期的结束；任务可以被抢占；任务之间相互独立。\ez \imp{周期越短（速率越高）的任务优先级越高}——“速率单调”的名字就是这么来的。课件写的是“执行时间较短的优先级高”，考试时按“周期”答更稳妥。RM 实现简单、行为可预测。\textbf{充分条件}：$\sum C_i/P_i \le n(2^{1/n}-1)$。$n=1$ 时上界为 1，$n=2$ 约为 0.828，$n=3$ 约为 0.780，$n \to \infty$ 时趋于 $\ln 2 \approx 0.693$。利用率低于上界一定可调度；介于上界和 1 之间则不一定，需要具体分析。

\textbf{最早截止时间优先 EDF}（Earliest Deadline First）：\textbf{动态优先级}算法，截止时间越早的任务优先级越高；每当有任务到达或完成时，重新比较各任务的截止时间。\textbf{优点}：对于可抢占的周期任务，在单处理器上只要 $\sum C_i/P_i \le 1$ 就\textbf{一定}可以调度，利用率可以达到 100\%，是单处理器上的最优算法。\textbf{缺点}：实现复杂，每次都要比较截止时间，开销大；需要知道精确的截止时间；系统过载时会出现\textbf{多米诺效应}，许多任务接连错过截止时间（而 RM 过载时至少能保证高优先级任务）。

例：T1（$C=2$，$P=5$），T2（$C=4$，$P=7$），利用率 $2/5 + 4/7 \approx 0.971$，超过两个任务的 RM 上界 0.828。\textbf{RM}：T1 周期短，优先级高。T1 运行 0–2，T2 运行 2–5；t=5 时 T1 的第二个周期到达，抢占 T2，运行 5–7；t=7 时 T2 的截止时间到了，它还差 1 个单位没完成，\textbf{错过截止时间}。\textbf{EDF}：t=5 时 T1 新周期的截止时间是 10，而 T2 的截止时间是 7，更早，所以 T2 继续运行，t=6 完成；T1 运行 6–8，满足截止时间 10。

Linux 的实时调度：\cd{SCHED_FIFO}、\cd{SCHED_RR}（实时优先级 1–99）和 \cd{SCHED_DEADLINE}（EDF + CBS 带宽限制）。

\subsection{实例：Linux 调度演化（CFS 细节不要求记忆）}

\textbf{进程分类与调度策略}：
\begin{itemize}
    \item \textbf{实时进程}：\cd{SCHED_FIFO}（同优先级先来先服务，没有时间片，一直运行到阻塞、主动让出或被更高优先级抢占）、\cd{SCHED_RR}（同优先级轮转，时间片 100ms）、\cd{SCHED_DEADLINE}。\imp{实时进程的优先级一定高于普通进程}，它们属于单独的调度类。
    \item \textbf{普通进程}：\cd{SCHED_NORMAL}（即 \cd{SCHED_OTHER}）、\cd{SCHED_BATCH}（批处理，不抢占交互进程）、\cd{SCHED_IDLE}（极低优先级），由 CFS/EEVDF 调度。
\end{itemize}
内核中的 \cd{prio} 取值 0–139，数值越小优先级越高：0–99 为实时进程，100–139 为普通进程；普通进程 \cd{static_prio = 120 + nice}。\ez \imp{nice 值的范围是 $-20$ 到 19，越小优先级越高}；普通用户只能调大 nice（降低自己的优先级），调小需要 root 权限。

\textbf{演化主线}（按“问题 $\to$ 改进 $\to$ 新问题”的链条记忆）：2.4 的 O(n) 调度器 $\to$ O(1) 调度器（Ingo Molnar）$\to$ SD 楼梯调度 $\to$ RSDL（Con Kolivas）$\to$ CFS（2.6.23，Molnar，O(log n)）$\to$ EEVDF（6.6 以后）。思路的变化：从“用经验公式猜测哪些是交互式进程并加以照顾”，到“不猜了，只追求公平”，再到“公平 + 延迟”。

\textbf{Linux 2.4 调度器}：
\begin{itemize}
    \item 只有一个全局的运行队列。每次调度都要\textbf{遍历所有就绪进程}，用 \cd{goodness()} 计算每个进程的“好坏程度”（普通进程约为 counter + 20 $-$ nice，实时进程为 1000 + 实时优先级），选最大的，所以是 O(n) 的。
    \item \cd{counter} 就是时间片，每次时钟中断减 1。当\textbf{所有}就绪进程的时间片都用完时，才统一重新计算一次（称为一个 epoch）。
    \item 重新计算时\textbf{连睡眠的进程也一起算}：\cd{counter = counter/2 + 基础值}。睡眠进程上一轮没用完的 counter 会叠加一半，所以经常睡眠的进程时间片更长、优先级更高，从而偏好 I/O 密集型进程。例如基础值为 6 时：CPU 密集型进程每轮都用完，重算后为 6；常睡眠的进程可达 9；一直睡眠的进程收敛到 12。
\end{itemize}
\bk \textbf{2.4 调度器的缺点}（老师明确说是考点）：(1) \textbf{可扩展性差}：每次调度都要遍历整个队列，开销随进程数增长；所有 CPU 共用一个全局队列和一把锁。(2) \textbf{高负载下性能差}：进程多时，重新计算时间片的开销也很大。(3) \textbf{用睡眠时间判断交互式进程不准确}：有些睡眠很多的进程并不是交互式的（如数据库后台进程）。(4) \textbf{内核不可抢占}：进程在内核态运行时不能被切换，高优先级的实时进程到来时也不能立即抢占，实时性差。

\textbf{O(1) 调度器}（2.6 早期）：
\begin{itemize}
    \item 动态优先级 $= \max(100, \min(\text{静态优先级} - \text{bonus} + 5, 139))$，其中 bonus 在 0–10 之间，由进程的\textbf{平均睡眠时间}决定：睡得越多 bonus 越大，优先级越高。
    \item 判定为交互式进程的条件：动态优先级 $\le 3 \times \text{静态优先级}/4 + 28$（对 nice 为 0 的进程，需要 bonus $\ge 7$）。
    \item 时间片：静态优先级小于 120 时为 $(140 - \text{静态优先级}) \times 20$ms，否则为 $(140 - \text{静态优先级}) \times 5$ms（nice 为 0 的进程为 100ms）。
    \item 数据结构：每个 CPU 有两个优先级数组 \textbf{active} 和 \textbf{expired}，每个数组有 140 个链表（每个优先级一个）和一个 140 位的\textbf{位图}，用“查找第一个置位位”指令在 O(1) 时间内找到最高优先级的非空链表。
    \item 时间片用完时：交互式进程和实时进程重置时间片后放回 active 数组（但如果 expired 中的进程已经等了太久，也会把它移入 expired，防止饥饿）；其他进程重新计算时间片后移入 expired 数组。active 数组为空时，\textbf{交换两个数组的指针}，这也是 O(1) 的，不需要像 2.4 那样遍历重算。
    \item \textbf{缺点}：判断交互式进程的经验公式复杂难懂，难以维护；误判仍然存在，交互响应仍不够好。
\end{itemize}

\textbf{SD 楼梯调度}：抛弃动态优先级，追求完全公平。把各个优先级看作楼梯的台阶，进程用完一个时间片就\textbf{下一级台阶}（进程本身的优先级不变，只是在这一轮中所处的级别降低）；下到最底层后，回到它初始优先级的下一级，并获得 (N+1) 倍的时间片。高优先级进程逐级下降，最终会与低优先级进程在同一层竞争，所以\textbf{不会饥饿}；交互式进程大部分时间在睡眠，别的进程在下楼梯时它还停在高处，醒来时仍有较高的优先级，于是\textbf{自然地照顾了交互式进程，不需要经验公式}。

\textbf{RSDL}（旋转楼梯截止时间调度）：重新引入 expired 数组；每个优先级有一个\textbf{组配额} $T_g$，组内每个进程有一个\textbf{进程配额} $T_p$（小于时间片）。进程用完 $T_p$ 就降一级（minor rotation）；\textbf{当某一级的组配额 $T_g$ 用完时，该级的所有进程都被强制降一级}，即使它们自己的 $T_p$ 还没用完——这相当于设置了截止时间，保证低优先级进程的等待时间是\textbf{可预测}的。进程用完整个时间片后放入 expired 数组的初始优先级处；active 为空或所有进程都降到了最底层时，交换 active 和 expired（major rotation）。RSDL 没有被合并进内核，但启发了 CFS。

\textbf{CFS 完全公平调度器的思想}：
\begin{itemize}
    \item 模拟一个“理想的多任务 CPU”：如果有 n 个进程，就让它们同时以 1/n 的速度运行。现实中 CPU 一次只能运行一个，CFS 就让“落后”最多的进程先运行。
    \item 不再跟踪睡眠时间，也不再区分交互式进程。每个进程维护一个\textbf{虚拟运行时间 vruntime}，它与实际运行时间成正比，与进程的权重成反比。调度时总是\textbf{选择 vruntime 最小的进程}。
    \item 用\textbf{红黑树}按 vruntime 组织就绪进程，选择的就是\textbf{最左边的节点}（缓存了这个节点，所以选取是 O(1)，插入和删除是 O(log n)）。正在运行的进程从树中移出，被切换下来时以新的 vruntime 重新插入（它的 vruntime 变大了，会插到偏右的位置）。所以树中的进程整体上从右向左移动。
    \item \textbf{公式}：在一个调度周期内，每个进程分配到的时间 = 调度周期 $\times$ 本进程权重 / 所有就绪进程的权重之和。\ez \imp{vruntime 增量 = 实际运行时间 $\times$ NICE\_0\_LOAD（1024）/ 本进程权重}。代入可知，在一个调度周期内各进程的 vruntime 增量都相同（权重被约掉了），这正是“公平”的含义。
    \item \textbf{权重}：nice 为 0 的权重是 1024，nice 每相差 1，权重约相差 1.25 倍，对应约 10\% 的 CPU 份额。nice $-20$ 的权重为 88761，nice 5 为 335，nice 19 为 15。
    \item 例：nice 0 和 nice 5 两个进程，调度周期 6ms。分配时间分别为 $6 \times 1024/1359 \approx 4.52$ms 和 $6 \times 335/1359 \approx 1.48$ms；两者的 vruntime 增量都是 4.52ms。实验测得 CPU 占用比例为 75.3\% : 24.6\%。
\end{itemize}
\textbf{课件中关于 CFS 的四个讨论问题}：
\begin{enumerate}
    \item \textbf{新进程的 vruntime 初值是不是 0？} 不是。如果是 0，新进程的 vruntime 远小于已运行很久的进程，会长期霸占 CPU。CFS 把它设为当前队列的 \cd{min_vruntime} 再加上一个调度周期中它应得的份额（新进程稍微靠后一点，防止通过不断 fork 来抢占 CPU）。
    \item \textbf{休眠进程的 vruntime 一直保持不变吗？} 休眠期间它不运行，vruntime 不增长；但唤醒时，CFS 把它设为 $\max(\text{原 vruntime}, \text{min\_vruntime} - \text{一个补偿值})$。这样它不能靠长时间睡眠积攒大量的“欠账”在醒来后长期霸占 CPU，但又能得到小幅照顾，偏向 I/O 密集型进程。
    \item \textbf{休眠进程被唤醒时会立刻抢占 CPU 吗？} 不一定。只有当它的 vruntime 比当前运行进程小，而且差值超过\textbf{唤醒粒度}时才抢占，以避免过于频繁的切换。
    \item \textbf{进程从一个 CPU 迁移到另一个 CPU 时 vruntime 会不会变？} 会。各 CPU 队列的 \cd{min_vruntime} 不同，所以迁出时减去原队列的 \cd{min_vruntime}（变成相对值），迁入时再加上新队列的 \cd{min_vruntime}。
\end{enumerate}
另外，fork 后可以配置让子进程先运行，因为子进程通常马上 exec，这样可以避免父进程先写内存而触发不必要的 COW 复制。

\textbf{CFS 的情景分析}：(1) \textbf{时钟中断}：更新当前进程的运行信息，把这一段的实际运行时间换算成 vruntime 加上去；如果它这一段运行的时间已超过本调度周期分配给它的时间，就设置需要调度的标记，随后进行一次调度；否则继续运行。(2) \textbf{主动调度}：更新当前进程的运行信息，把它插回红黑树，选出最左节点，把它从树中移出并运行。(3) \textbf{进程唤醒}：调整其 vruntime 后插入红黑树，判断是否抢占当前进程。(4) \textbf{进程创建}：设置 vruntime 初值并插入红黑树。

\textbf{CFS 的数据结构}：\cd{task_struct} 中嵌入 \cd{struct sched_entity se}（包含权重 \cd{load}、红黑树节点 \cd{run_node}、\cd{vruntime}、累计运行时间 \cd{sum_exec_runtime}），它包含调度单个任务或任务组所需的完整信息；\cd{struct sched_class} 是调度类的接口，每个调度模块实现其中的一组函数（体现机制与策略分离）；每个 CPU 的运行队列中有一个 \cd{struct cfs_rq}，保存红黑树的根和 \cd{min_vruntime}。

\textbf{EEVDF}（课堂派作业）：在 CFS“公平”的基础上引入了\textbf{延迟}这一因素。定义 \textbf{lag} = 应得的服务 $-$ 实际得到的服务；lag $\ge 0$ 的进程是 \textbf{eligible}（有资格运行的）。每个进程有一个\textbf{虚拟截止时间} = eligible 时间 + 请求的时间片 / 权重。调度时\textbf{在所有 eligible 的进程中选虚拟截止时间最早的}。请求较短时间片的进程虚拟截止时间早，能更快得到运行（延迟低），但每次只运行一小段；它的总份额仍由权重决定。这样就把\textbf{延迟和份额解耦}了，可以看作公平共享与 EDF 的结合。

\subsection{xv6 源码阅读：调度（swtch.S、proc.h、proc.c）}

\textbf{数据结构}：\cd{struct context} 保存 \cd{ra}、\cd{sp}、\cd{s0}–\cd{s11} 共 14 个寄存器；\cd{struct cpu} 包含 \cd{proc}（当前在这个 CPU 上运行的进程）、\cd{context}（该 CPU 调度器线程的上下文）、\cd{noff}（\cd{push_off} 关中断的嵌套深度）、\cd{intena}（第一次 \cd{push_off} 之前中断是否打开）。\cd{struct proc} 中与调度相关的\textbf{只有 \cd{state}}，没有优先级、时间片之类的字段。

\textbf{\cd{scheduler()}}：每个 CPU 在 \cd{main} 的最后调用它，永不返回。外层是 \cd{for(;;)} 无限循环：先 \cd{intr_on()} 开中断（避免所有进程都在等待时死锁——必须让设备中断能进来唤醒它们）；内层扫描整个 \cd{proc[]} 数组，对每个进程先 \cd{acquire(&p->lock)}，如果它是 RUNNABLE，就设置 \cd{p->state = RUNNING; c->proc = p;}，然后 \cd{swtch(&c->context, &p->context)} 切换到该进程；该进程放弃 CPU 后，执行流从这个 \cd{swtch} 返回，设置 \cd{c->proc = 0}；最后 \cd{release(&p->lock)}。如果扫描一整轮都没有找到可运行的进程，就执行 \cd{wfi} 等待中断。调度策略是\textbf{按进程表顺序轮转、没有优先级}，时间片就是一个时钟中断间隔（每次时钟中断都 \cd{yield}）。

\textbf{\cd{yield()}}：依次执行 \cd{acquire(&p->lock);} \cd{p->state = RUNNABLE;} \cd{sched();} \cd{release(&p->lock);}。

\textbf{\cd{sched()}}：切换前做四项检查，不满足就 panic：(1) 当前进程持有 \cd{p->lock}；(2) \cd{mycpu()->noff == 1}，即只持有这一把锁（如果还持有别的锁就切走，其他进程可能因为等这把锁而死锁）；(3) 进程状态已经不是 RUNNING（调用者应该已经改好了）；(4) 中断已关闭。然后保存 \cd{intena}（它属于这个内核线程，而不属于 CPU），调用 \cd{swtch(&p->context, &mycpu()->context)} 切换到调度器线程，回来后恢复 \cd{intena}。调用 \cd{sched()} 的地方有三处：\cd{yield}（时间片用完）、\cd{sleep}（阻塞）、\cd{exit}（退出，永不返回）。

\textbf{一次从进程 A 切换到进程 B 的完整过程}：A 在用户态发生时钟中断 $\to$ \cd{usertrap} 调用 \cd{yield}，获取 A 的锁，把 A 设为 RUNNABLE $\to$ \cd{sched} $\to$ \cd{swtch} 切换到本 CPU 的调度器线程 $\to$ 调度器从它上次调用的 \cd{swtch} 中返回，\textbf{释放 A 的锁} $\to$ 继续扫描，找到 B，获取 B 的锁，把 B 设为 RUNNING $\to$ \cd{swtch} 切换到 B $\to$ B 从它上次在 \cd{sched} 中调用的 \cd{swtch} 返回，回到 \cd{yield}，\textbf{释放 B 的锁} $\to$ \cd{usertrapret} $\to$ B 回到用户态。

要点：
\begin{enumerate}
    \item \textbf{一次进程切换要执行两次 \cd{swtch}}（A $\to$ 调度器 $\to$ B），\cd{sched} 和 \cd{scheduler} 互相切换，像一对协程。使用独立的调度器线程和栈，使“选下一个进程”这件事与任何进程都无关：A 可能正在执行 \cd{exit}，它的内核栈马上要被回收，不能在 A 的栈上运行调度代码；调度器线程还为“没有进程可运行”时提供了一个待着的地方。Linux 则直接在当前进程的内核栈上调用 \cd{schedule()}，只切换一次。
    \item \cd{p->lock} \textbf{跨越 \cd{swtch} 传递}：在 A 中获取的锁，在调度器中释放；在调度器中获取的 B 的锁，在 B 中释放。这把锁保护的是“进程状态与其 \cd{context} 一致”这一条件。如果 A 在调用 \cd{swtch} 之前就释放了锁，另一个 CPU 的调度器可能看到 A 已经是 RUNNABLE，立刻切换到 A，而此时 A 的寄存器还没有保存完——\textbf{两个 CPU 会同时使用 A 的内核栈}。持有锁的同时中断也是关闭的，防止切换到一半又因时钟中断而再次 \cd{yield}。
\end{enumerate}

\textbf{新进程第一次被调度}：\cd{allocproc} 已把 \cd{context.ra} 设为 \cd{forkret}，\cd{context.sp} 设为该进程内核栈的顶部。调度器 \cd{swtch} 到新进程后，\cd{ret} 跳到 \cd{forkret}：它释放 \cd{p->lock}（与上面锁的传递相呼应），如果是系统中的第一个进程，还要调用 \cd{fsinit} 初始化文件系统（这一步可能需要 \cd{sleep} 等磁盘，所以必须在进程上下文中完成，不能放在 \cd{main} 里）；然后调用 \cd{usertrapret} 返回用户态，于是子进程看到 fork 返回 0。

\textbf{\cd{mycpu()} 与 \cd{myproc()}}：\cd{mycpu()} 读取 \cd{tp} 寄存器（存放 hartid）来确定当前 CPU，必须在关中断的情况下使用，否则读完 \cd{tp} 后可能被时钟中断打断并迁移到另一个 CPU 上，拿到的 \cd{cpu} 结构就不对了。\cd{myproc()} 内部自己做了 \cd{push_off}/\cd{pop_off}，返回的进程指针在之后即使发生迁移也仍然有效。

\textbf{xv6-k210 的差异}：每个进程有自己的内核页表，所以切换前后要切换 \cd{satp} 并刷新 TLB；时钟中断通过 SBI 的 \cd{sbi_set_timer} 设置；\cd{timer_tick} 中统计进程的 utime/stime；\cd{forkret} 中初始化的是 FAT32 文件系统；\cd{tp} 在 \cd{main} 中由 SBI 传入的 hartid 设置。
